\documentclass[../../main.tex]{subfiles}
\begin{document}
\section{余向量场的拉回}\label{section:ch11-余向量场的拉回}

正如我们所见，光滑映射在切向量上给出了一个称为微分的线性映射。对偶化这一映射，就得到一个方向相反的、在余向量上的线性映射。

设 $F\colon M\to N$ 是带边或不带边光滑流形间的光滑映射，$p\in M$ 任意。微分 $dF_p\colon T_pM\to T_{F(p)}N$ 给出对偶线性映射
\begin{align*}
dF_p^*\colon T_{F(p)}^*N\to T_p^*M,
\end{align*}
称为 $F$ 在 $p$ 处的\textbf{（逐点）拉回}，或 $F$ 的\textbf{余切映射}。展开定义，可见 $dF_p^*$ 的特征是
\begin{align}
dF_p^*(\omega)(v)=\omega(dF_p(v)),\quad \omega\in T_{F(p)}^*N,\ v\in T_pM.\label{align:ch11-逐点拉回}
\end{align}

注意到指派 $(M,p)\mapsto T_p^*M$ 与 $F\mapsto dF_p^*$ 给出了从带点光滑流形范畴到实向量空间范畴的一个反变函子。正因如此，把 $T^*M$ 的元素称为“协变向量”的约定尤为不幸；但这一术语如此根深蒂固，我们别无选择只能沿用。

在余向量场上，逐点拉回结合为余向量场的拉回。若 $F\colon M\to N$ 是光滑映射，$\omega$ 是 $N$ 上的光滑余向量场，则 $F^*\omega$ 是 $M$ 上的余向量场，定义为
\begin{align}
(F^*\omega)_p=dF_p^*(\omega_{F(p)}).\label{align:ch11-余向量场拉回2}
\end{align}

\textbf{例 11.28}。若 $F\colon M\to N$ 是光滑映射，而 $f\in C^\infty(N)$，则 $f$ 的微分被拉回为
\begin{align*}
(F^*(df))_p(v)=df_{F(p)}(dF_p(v))=dF_p(v)(f)=v(f\circ F)=d(f\circ F)_p(v).
\end{align*}
于是
\begin{align}
d(f\circ F)=F^*(df).\label{align:ch11-链法则微分}
\end{align}
这只是链式法则的余向量形式。$\blacktriangleright$

下面的命题给出了在坐标下计算拉回的公式。

\begin{proposition}[拉回的坐标公式]\label{proposition:ch11-拉回坐标公式}
设 $F\colon M\to N$ 是带边或不带边光滑流形间的光滑映射，$(y^j)$ 是 $N$ 上的局部坐标，$(x^i)$ 是 $M$ 上的局部坐标。若 $\omega=\omega_j\,dy^j$ 是 $N$ 上的光滑余向量场，则
\begin{align}
F^*\omega=\omega_j\circ F\,\frac{\partial F^j}{\partial x^i}\,dx^i.\label{align:ch11-拉回坐标公式}
\end{align}
\end{proposition}

\begin{proof}
由\eqref{align:ch11-链法则微分}，$F^*(dy^j)=d(y^j\circ F)=dF^j=(\partial F^j/\partial x^i)dx^i$。因此由拉回的线性性，
\begin{align*}
F^*\omega=F^*(\omega_j\,dy^j)=(\omega_j\circ F)F^*(dy^j)=\omega_j\circ F\,\frac{\partial F^j}{\partial x^i}\,dx^i.
\end{align*}
\end{proof}

与切向量场形成鲜明对比，余向量场在光滑映射下总能被拉回（前提是映射方向正确）。这意味着拉回给出了余向量场的一个反变函子。

\begin{proposition}[拉回的函子性质]\label{proposition:ch11-拉回函子}
设 $F\colon M\to N$ 与 $G\colon N\to P$ 为光滑流形间的光滑映射。则对任意 $\omega\in\mathfrak X^*(P)$，
\begin{align}
(G\circ F)^*\omega=F^*(G^*\omega).\label{align:ch11-拉回复合}
\end{align}
此外，对 $M$ 上的恒等映射 $\mathrm{Id}_M$，有
\begin{align}
\mathrm{Id}_M^*\omega=\omega,\quad \omega\in\mathfrak X^*(M).\label{align:ch11-恒等拉回}
\end{align}
\end{proposition}

\begin{proof}
逐点拉回的相应性质，即对偶函子的反变性，是线性代数中的标准事实。
\end{proof}

下面的命题表明，余向量场的拉回与限制到子流形是相容的。

\begin{proposition}[拉回与限制]\label{proposition:ch11-拉回与限制}
设 $S\subseteq M$ 是带边或不带边的光滑子流形，$\iota\colon S\hookrightarrow M$ 为包含映射。则对 $M$ 上的任意光滑余向量场 $\omega$，拉回 $\iota^*\omega$ 恰为 $\omega$ 限制到 $S$ 上所得的余向量场。
\end{proposition}

\begin{proof}
这直接由定义得到：对 $p\in S$ 与 $v\in T_pS$，$d\iota_p(v)=v$（视为 $T_pM$ 中的元素），故
\begin{align*}
(\iota^*\omega)_p(v)=\omega_p(d\iota_p(v))=\omega_p(v).
\end{align*}
\end{proof}
\end{document}
